\documentclass{article}

\usepackage{amsmath, amsthm, amssymb, amsfonts}
\usepackage{thmtools}
\usepackage{setspace}
\usepackage{geometry}
\usepackage[utf8]{inputenc}
\usepackage[english]{babel}
\usepackage{framed}
\usepackage[dvipsnames]{xcolor}
\usepackage{tcolorbox}

\colorlet{LightGray}{White!90!Periwinkle}
\colorlet{LightOrange}{Orange!15}
\colorlet{LightGreen}{Green!15}
\definecolor{LightBlue}{RGB}{137, 207, 240}

\newcommand{\HRule}[1]{\rule{\linewidth}{#1}}
\newcommand{\norm}[1]{\left \lVert #1 \right \rVert}

\declaretheoremstyle[name=Theorem,]{thmsty}
\declaretheorem[style=thmsty]{theorem}
\tcolorboxenvironment{theorem}{colback=LightGray}

\declaretheoremstyle[name=Proposition,]{prosty}
\declaretheorem[style=prosty,numberlike=theorem]{proposition}
\tcolorboxenvironment{proposition}{colback=LightOrange}

\declaretheoremstyle[name=Principle,]{prcpsty}
\declaretheorem[style=prcpsty,numberlike=theorem]{principle}
\tcolorboxenvironment{principle}{colback=LightGreen}

\declaretheoremstyle[name=Definition,]{defsty}
\declaretheorem[style=defsty]{definition}
\tcolorboxenvironment{definition}{colback=LightBlue!50}

\setstretch{1.2}
\geometry{
    textheight=9in,
    textwidth=5.5in,
    top=1in,
    headheight=12pt,
    headsep=25pt,
    footskip=30pt
}

% ------------------------------------------------------------------------------

\begin{document}

% ------------------------------------------------------------------------------
% Cover Page and ToC
% ------------------------------------------------------------------------------

\title{ \normalsize \textsc{}
\\ [2.0cm]
\HRule{1.5pt} \\
\LARGE \textbf{\uppercase{Math 291}
\HRule{2.0pt} \\ [0.6cm] \LARGE{Honors Multivariable Calculus} \vspace*{10\baselineskip}}
}
\date{}
\author{\textbf{Jeffrey Lu}}

\maketitle
\newpage

\tableofcontents
\newpage

% ------------------------------------------------------------------------------

\section{Lines and Planes}

% Theorem 1
\begin{theorem}
    Let $V \subset \mathbb{R}^n,$ and let $W = Span(V)$. Then $Span(W) = W$
\end{theorem}

% Theorem 2
\begin{theorem}
    (Fundamental property of the standard basis).
    Let ${\mathbf{e_1},...,\mathbf{e_n}}$ be the standard basis in $\mathbb{R}^n$.
    $Span({\mathbf{e_1},...,\mathbf{e_n}}) = \mathbb{R}^n$,
    and for each $\mathbf{x} \in \mathbb{R}^n$,
    there is exactly one vector of coefficients $(t_1,...,t_n)$ such that
    $$\mathbf{x} = \sum_{j=1}^nt_j\mathbf{e}_j$$
    namely $(t_1,...,t_n) = x_1,...,x_n)$ .
\end{theorem}

% Theorem 3
\begin{theorem}
    (Cauchy–Schwartz inequality).
    For any two vectors $\mathbf{a, b} \in \mathbb{R}^n$,
    $$|\mathbf{a}\cdot\mathbf{b}|\le \mathbf{\norm{a}\norm{b}}$$
    There is a inequality iff $\norm{\mathbf{b}}\mathbf{a}=\pm\mathbf{\norm{a}b}$
\end{theorem}

% Theorem 4
\begin{theorem}
    (Triangle inequality).
    For any three vectors $\mathbf{x, y, z} \in \mathbb{R}^n$,
    $$\norm{\mathbf{x-z}} \le \norm{\mathbf{x-z}} + \norm{\mathbf{y-z}}$$
\end{theorem}

% Theorem 5
\begin{theorem}
    The Euclidean distance function
    $$\varrho_E(x,y):=\mathbf{\sqrt{(x-y)\cdot(x-y}=\norm{x-y}}$$
    is a metric in $\mathbb{R}^n$
\end{theorem}
% Theorem 6
\begin{theorem}
    (Fundamental Theorem on Orthonormal Sets in $\mathbb{R}^n$).
    Let $\{\mathbf{u_1},...,\mathbf{u_n}\}$ be any orthonormal set in $\mathbb{R}^n$ consisting of exactly $n$ vectors. Then every vector $\mathbf{x}\in\mathbb{R}^n$ can be written as a linear combination of the vectors $\{\mathbf{u_1},...,\mathbf{u_n}\}$ in exactly one way, namely
    $$\mathbf{x}=\sum_{j=1}^n(\mathbf{x}\cdot\mathbf{u_j})\mathbf{u_j}$$
    The squared length of $\mathbf{x}$ is the sum of the squares of the coefficients in the expansion:
    $$\mathbf{\norm{x}}=\sum_{j=1}^n(\mathbf{x}\cdot\mathbf{u_j})^2$$
    The standard basis of $\mathbb{R}^n, \{\mathbf{e_1}, . . . , \mathbf{e_n}\}$ is a set of $n$ orthonormal vectors in $\mathbb{R}^n$, and so the theorem says that
    $$\mathbf{x}=\sum_{j=1}^nx_j\mathbf{e_j}=\sum_{j-1}^n(\mathbf{x}\cdot\mathbf{e_j})\mathbf{e_j}$$
    is the unique way to express any $\mathbf{x}$ in $\mathbb{R}^n$ as a linear combination of the standard basis vectors.

\end{theorem}

% Definition 1
\begin{definition}
    (Linear Combinations and Span). Let $V \subset \mathbb{R}$. The span of $V$ is the set of all possible linear combinations $\sum_{j=1}^mt_j\mathbf{v_j}$.
\end{definition}

% Theorem 7
\begin{theorem}
    (Algebraic properties of the cross product).
    Let $\mathbf{a, b,}$ and $\mathbf{c}$ be any three vectors in $\mathbb{R}^3$, and let $t$ be any number. Then
    \begin{gather*}
        \mathbf{a\times{b}=-b\times{a}} \\
        \mathbf{a\times(a+b)=a\times{b}+a\times{c}} \\
        (t\mathbf{a})\times{\mathbf{b}} = t(\mathbf{a\times{b})}= \mathbf{a}\times{(t\mathbf{b})}
    \end{gather*}

\end{theorem}

% Theorem 8
\begin{theorem}
    (Triple product identity).
    Let $\mathbf{a, b, c} \in \mathbb{R}^n$. Then
    $$(\mathbf{a\times{b})\cdot{c}=(b\times{c})\cdot{a}}$$
    In other words, the triple product $\mathbf{(a\times{b})\cdot{c}}$ is unchanged when the vectors in it are cyclicly permuted.
\end{theorem}

% Theorem 9
\begin{theorem}
    (Cross product equals normal vector)
    Let $\mathbf{a, b}\in\mathbb{R}^3.$ Then $\mathbf{a\times{b}}$ is orthogonal to both $\mathbf{a}$ and $\mathbf{b}$.
\end{theorem}

% Theorem 10
\begin{theorem}
    Let $\{\mathbf{u_1, u_2, u_3}\}$ be any orthonormal basis of $\mathbb{R}^3$. Then
    $$\mathbf{u_1\times{u_2}=u_3}\Longleftrightarrow\mathbf{u_2\times{u_3}=u_3}\Longleftrightarrow\mathbf{u_3\times{u_1}=u_2}$$
    In particular, $\{\mathbf{u_1, u_2, u_3}\}$ is right handed if and only if any one of the identities in equation above is valid,
    and in that case, all of them are valid.
\end{theorem}

% Theorem 11
\begin{theorem}
    (Invariance under change of basis).
    Let $\{\mathbf{u_1, u_2, u_3}\}$ be any right handed orthonormal
    basis in $\mathbb{R}^3$. For any vectors $\mathbf{x} = (x_1, x_2, x_3)$ and $\mathbf{y} = (y_1, y_2, y_3)$, define $\mathbf{a} = x_1\mathbf{u_1} + x_2\mathbf{u_2} + x_3\mathbf{u_3}$and $\mathbf{b} = y_1\mathbf{u_1} + y_2\mathbf{u_2} + y_3\mathbf{u_3}$, so that $\mathbf{x}$ and $\mathbf{y}$ are the coordinate vectors of $\mathbf{a}$ and $\mathbf{b}$ with respect to the coordinate system based on the orthonormal basis $\{\mathbf{u_1, u_2, u_3}\}$. Then
    $$\mathbf{a\times{b}}=(x_2y_3-x_3y_2)\mathbf{u_1}+(x_3y_1-x_1y_3)\mathbf{u_2}+(x_1y_2-x_2y_1)\mathbf{u_3}$$
    In other words, one can compute the cross product of $\mathbf{a}$ and $\mathbf{b}$ by computing the cross product of their
    coordinate vectors for any right handed orthonormal basis.
\end{theorem}

% Theorem 12
\begin{theorem}
    (Lagrange's Identity.) Let $\mathbf{a, b, c} \in \mathbb{R}$. Then $$\mathbf{a} \times (\mathbf{b} \times \mathbf{c})=(\mathbf{a} \cdot \mathbf{c})\mathbf{b} - (\mathbf{a} \cdot \mathbf{b})\mathbf{c}$$
\end{theorem}
how fast

\end{document}
